% AIビジネス入門 第09回 AIを構成する技術(1)：記号とルール（CC BY-SA 4.0）
\documentclass[aspectratio=169]{beamer}
\usepackage{lfs-slides}
\usepackage{url}
\input{../aib-common.tex.inc}
\lfscourse{AIビジネス入門}

\title[AIを構成する技術(1)：記号とルール]{AIを構成する技術(1)\\記号とルール}
\subtitle{AIビジネス入門 第09回}
\author{清田 陽司}
\date{}

\begin{document}

\begin{frame}
\titlepage
\end{frame}

% ---- 位置づけ ----
\begin{frame}{AI は「どう動くか」：2つの考え方}
\begin{center}
\begin{tikzpicture}[box/.style={draw,rounded corners,text width=4.6cm,minimum height=1.9cm,align=center,font=\small}]
  \node[box,fill=boom!10,draw=boom] (x) {\textbf{今日 第09回}\\{\large 記号とルール}\\人が知識を\textbf{書く}};
  \node[box,fill=winter!10,draw=winter,right=0.6cm of x] (y) {\textbf{次回 第10回}\\{\large 機械学習}\\データから\textbf{学ぶ}};
\end{tikzpicture}
\end{center}
\vspace{0.8em}
\centering 数式は使いません。\textbf{考え方の違い}をつかむ
\end{frame}

% ---- 今回のトピックス ----
\begin{frame}{今回のトピックス}
\begin{center}
\begin{tikzpicture}[box/.style={draw,rounded corners,text width=2.9cm,minimum height=2.8cm,align=center,font=\small}]
  \node[box,fill=gray!10] (a) {\textbf{AI の\\二大潮流}\\[0.5em]記号主義と\\コネクショニズム};
  \node[box,fill=gain!8,right=0.3cm of a] (b) {\textbf{知識表現}\\[0.5em]ルール\\分類・つながり\\アフォーダンス};
  \node[box,fill=winter!8,right=0.3cm of b] (c) {\textbf{強みと限界}\\[0.5em]説明できる\\でも\\書ききれない};
  \node[box,fill=boom!8,right=0.3cm of c] (d) {\textbf{ビジネスと\\倫理}\\[0.5em]融資・医療・流通\\だれのルールか};
\end{tikzpicture}\\[0.5em]
\begin{tikzpicture}
  \node[draw,rounded corners,fill=genai!12,draw=genai,text width=13.2cm,align=center,inner sep=7pt] {\textbf{生成 AI は、記号的 AI の限界を超えたのか？} \quad ハルシネーションと RAG};
\end{tikzpicture}
\end{center}
\end{frame}

% ============================================================
% 二大潮流
% ============================================================

\sectionframe{AI の二大潮流\\[0.3em]記号主義とコネクショニズム}

\begin{frame}{ネコの写真を見分ける：2つのやり方}
\begin{center}
\begin{tikzpicture}[box/.style={draw,rounded corners,text width=5.8cm,minimum height=4.0cm,align=left,font=\small,inner sep=7pt}]
  \node[box,fill=boom!8,draw=boom] (a) {\centering\textbf{\large 記号主義}\\{\footnotesize 人が\textbf{ルールを書く}}\par\vspace{0.5em}
  \texttt{IF}\ 耳が三角形\\
  \quad \texttt{AND} ひげがある\\
  \quad \texttt{AND} しっぽが長い\\
  \texttt{THEN}「ネコである」\par\vspace{0.4em}
  {\footnotesize なぜそう判定したか、\textbf{読めば分かる}}};
  \node[box,fill=winter!8,draw=winter,right=0.6cm of a] (b) {\centering\textbf{\large コネクショニズム}\\{\footnotesize データから\textbf{学ばせる}}\par\vspace{0.5em}
  ネコの写真 1 万枚と\\ネコでない写真 1 万枚\\→ ちがいを\textbf{自動で}見つける\\→ パターンとして判定\par\vspace{0.4em}
  {\footnotesize ルールは書かない。\textbf{中身は読みにくい}}};
\end{tikzpicture}
\end{center}
\vspace{0.3em}
\centering{\large ほとんど\textbf{正反対}のアプローチ}\quad\small 今日は左。右は次回
\end{frame}

\begin{frame}{2つの流れの歴史}
\begin{center}
\begin{tikzpicture}[x=1.45cm]
  \draw[-{Latex[length=3mm]},thick] (0,0) -- (8.6,0) node[right,font=\small] {年};
  \foreach \x/\y in {0.5/1960,1.5/1970,2.5/1980,3.5/1990,4.5/2000,5.5/2010,6.5/2020}{\draw (\x,0.08) -- (\x,-0.08); \node[below,font=\scriptsize] at (\x,-0.08) {\y};}
  \fill[boom!25] (0.3,-0.75) rectangle (1.9,-1.15); \node[font=\scriptsize\bfseries] at (1.1,-0.95) {第1次ブーム};
  \fill[boom!25] (2.6,-0.75) rectangle (3.7,-1.15); \node[font=\scriptsize\bfseries] at (3.15,-0.95) {第2次};
  \fill[boom!25] (5.6,-0.75) rectangle (6.6,-1.15); \node[font=\scriptsize\bfseries] at (6.1,-0.95) {第3次};
  \fill[genai!30] (6.75,-0.75) rectangle (8.3,-1.15); \node[font=\scriptsize\bfseries] at (7.5,-0.95) {生成 AI};
  \draw[boom,line width=3pt,-{Latex[length=3mm]}] (0.2,1.0) -- (8.3,1.0);
  \node[anchor=west,font=\small\bfseries,boom!80!black] at (0.1,1.45) {記号主義};
  \foreach \x/\t in {1.0/推論・探索,3.1/エキスパートシステム,5.4/Web とナレッジグラフ}{\node[draw=boom,fill=white,rounded corners,font=\scriptsize] at (\x,1.0) {\t};}
  \draw[winter,line width=3pt,-{Latex[length=3mm]}] (0.2,2.3) -- (8.3,2.3);
  \node[anchor=west,font=\small\bfseries,winter!80!black] at (0.1,2.75) {コネクショニズム};
  \foreach \x/\t in {0.9/ニューラルネット誕生,3.2/誤差逆伝播法,6.0/ディープラーニング}{\node[draw=winter,fill=white,rounded corners,font=\scriptsize] at (\x,2.3) {\t};}
\end{tikzpicture}
\end{center}
\vspace{0.4em}
\begin{itemize}\small
  \item 長いあいだ AI の\textbf{主流は記号主義}。\textbf{第2次ブーム}（エキスパートシステム）がその頂点
  \item 2010年代からはコネクショニズムが全盛。でも記号主義は\textbf{消えていない}（今日の後半）
\end{itemize}
{\scriptsize 年代は大づかみの目安}
\end{frame}

\thinkframe{みなさんの毎日の行動で\\[0.3em]「\textbf{もし〜なら、〜する}」\\[0.3em]で動いているものを1つ\\[0.6em]{\large 例：もし残高が 1,000 円を切ったら、チャージする}}

% ============================================================
% 知識表現
% ============================================================

\sectionframe{記号的 AI の基礎\\[0.3em]知識表現}

\begin{frame}{知識表現：世界と AI のあいだの「通訳」}
\begin{center}
\begin{tikzpicture}[box/.style={draw,rounded corners,minimum height=1.5cm,align=center,font=\small},arr/.style={-{Latex[length=3mm]},very thick,gray}]
  \node[box,fill=gain!12,draw=gain,text width=3.2cm] (w) {\textbf{\large 現実の世界}\\机・りんご・車\\できごと};
  \node[box,fill=boom!15,draw=boom,line width=1.2pt,text width=3.6cm,right=0.9cm of w] (k) {\textbf{\large 知識表現}\\ことばと\\ことばの\textbf{関係}};
  \node[box,fill=winter!12,draw=winter,text width=3.2cm,right=0.9cm of k] (a) {\textbf{\large AI}\\判断する\\行動を決める};
  \draw[arr] (w) -- (k); \draw[arr] (k) -- (a);
\end{tikzpicture}\\[1.0em]
\begin{tikzpicture}[h/.style={draw,rounded corners,text width=5.6cm,minimum height=0.8cm,align=center,font=\small\bfseries},b/.style={draw,rounded corners,fill=gray!6,text width=5.6cm,minimum height=2.0cm,align=left,font=\small,inner sep=6pt,anchor=north}]
  \node[h,fill=gray!15] (h1) at (0,0) {AI には体がない};
  \node[b] at ($(h1.south)+(0,-0.1)$) {・人は見て、さわって、自然に分かる\\・AI は世界を\textbf{直接}は分からない\\・「どう見るか」の\textbf{枠組み}が要る};
  \node[h,fill=boom!12,draw=boom] (h2) at (6.4,0) {知識表現の役割};
  \node[b] at ($(h2.south)+(0,-0.1)$) {・世界の\textbf{見方}を与える\\・何を区別するかを決める\\・コンピュータが\textbf{処理できる形}に直す};
\end{tikzpicture}
\end{center}
\end{frame}

\begin{frame}{いちばん基本の形：「もし〜なら、〜する」}
\begin{columns}[c]
\begin{column}{0.46\textwidth}
\centering
{\Large\texttt{IF} 条件 \texttt{THEN} 行動}\\[1.2em]
\begin{tikzpicture}[box/.style={draw,rounded corners,text width=5.0cm,align=left,font=\small,inner sep=6pt}]
  \node[box,fill=boom!8,draw=boom] {\textbf{通販サイトのチャット窓口}\\[0.2em]
  ・\texttt{IF}「返品」と書かれた\\\quad \texttt{THEN} 返品の手順を案内する\\
  ・\texttt{IF} 注文から 30 日を過ぎた\\\quad \texttt{THEN}「返品期限を過ぎています」\\
  ・\texttt{IF} どれにも当てはまらない\\\quad \texttt{THEN} 担当者につなぐ};
\end{tikzpicture}
\end{column}
\begin{column}{0.52\textwidth}
\begin{tikzpicture}[box/.style={draw,rounded corners,text width=5.8cm,align=left,font=\small,inner sep=6pt}]
  \node[box,fill=gray!8] (a) {\textbf{ルールを\textbf{たくさん}用意すると}\\場面に合った案内ができる\\→ 「かしこく」見える};
  \node[box,fill=gain!10,draw=gain,below=0.3cm of a] (b) {\textbf{みなさんの例も同じ形}\\「もし 8 時を過ぎていたら、\\\quad 朝ごはんはトーストだけ」};
  \node[box,fill=winter!10,draw=winter,below=0.3cm of b] (c) {\textbf{だれが書く？}\\ルールは\textbf{人が書く}。\\AI が自分で作るのではない};
\end{tikzpicture}
\end{column}
\end{columns}
\end{frame}

\begin{frame}{ルールをつなぐ：ゴールから逆にたどる}
\begin{center}
\begin{tikzpicture}[box/.style={draw,rounded corners,minimum height=1.1cm,align=center,font=\small,text width=2.5cm},arr/.style={-{Latex[length=3mm]},very thick,gray}]
  \node[box,fill=boom!15,draw=boom,line width=1.2pt] (g) at (0,0) {\textbf{ゴール}\\発表会を開く};
  \node[box,fill=gray!8] (i) at (4,0.9) {会場の予約};
  \node[box,fill=winter!12,draw=winter] (t) at (4,-0.9) {招待状};
  \node[box,fill=gray!8] (p) at (8,0) {参加者の名簿};
  \node[box,fill=gray!8] (w) at (8,-1.8) {日程};
  \draw[arr] (i) -- (g); \draw[arr] (t) -- (g);
  \draw[arr] (p) -- (t); \draw[arr] (w) -- (t);
  \node[font=\scriptsize,text width=3.4cm,align=center] at (2.0,-2.0) {知識：発表会には\\会場と招待状が要る};
  \node[font=\scriptsize,text width=3.4cm,align=center] at (6.0,-2.6) {知識：招待状には\\名簿と日程が要る};
\end{tikzpicture}
\end{center}
\vspace{0.5em}
\centering ゴールを与えると、AI は\textbf{知識をつないで}「名簿と日程を決める → 招待状 → \ldots」と\textbf{計画}を立てる\\[0.3em]
\small 仕事の段取り（何をそろえれば何ができるか）を、ルールの組み合わせで表す
\end{frame}

\begin{frame}{オントロジー：ものごとの「分類」の体系}
\begin{center}
\begin{tikzpicture}[n/.style={draw,rounded corners,align=center,font=\small,minimum height=0.75cm,inner sep=4pt},l/.style={draw,rounded corners,align=center,font=\scriptsize,fill=gray!6,minimum height=0.6cm,inner sep=3pt},e/.style={gray,thick}]
  \node[n,fill=boom!15,draw=boom] (r) at (5,0) {\textbf{家電}};
  \node[n,fill=gain!12,draw=gain] (h) at (1.2,-1.3) {調理家電};
  \node[n,fill=winter!12,draw=winter] (a) at (5,-1.3) {生活家電};
  \node[n,fill=genai!12,draw=genai] (t) at (8.8,-1.3) {映像・音響};
  \foreach \x/\t in {-0.4/電子レンジ,1.2/炊飯器,2.6/トースター}{\node[l] (h\t) at (\x,-2.6) {\t}; \draw[e] (h.south) -- (h\t.north);}
  \foreach \x/\t in {3.9/掃除機,5/洗濯機,6.1/エアコン}{\node[l] (a\t) at (\x,-2.6) {\t}; \draw[e] (a.south) -- (a\t.north);}
  \foreach \x/\t in {8.2/テレビ,9.5/イヤホン}{\node[l] (t\t) at (\x,-2.6) {\t}; \draw[e] (t.south) -- (t\t.north);}
  \draw[e] (r.south) -- (h.north); \draw[e] (r.south) -- (a.north); \draw[e] (r.south) -- (t.north);
\end{tikzpicture}
\end{center}
\vspace{0.5em}
\begin{itemize}\small
  \item 「何が何の\textbf{仲間}か」を上下の関係で書く。AI は「この人は\textbf{調理家電}をよく見ている」→ 同じ仲間をすすめる、と使える
  \item ビジネスでは：\textbf{商品の分類}（EC のカテゴリ）、病名の分類、部品表
\end{itemize}
\end{frame}

\begin{frame}{セマンティック・ネットワーク：ことば同士の「つながり」}
\begin{columns}[c]
\begin{column}{0.56\textwidth}
\centering
\begin{tikzpicture}[n/.style={draw,rounded corners,align=center,font=\small,minimum height=0.8cm,minimum width=1.9cm},arr/.style={-{Latex[length=2.5mm]},thick},lab/.style={font=\scriptsize,fill=white,inner sep=1pt}]
  \node[n,fill=boom!15,draw=boom] (bird) at (0,0) {\textbf{鳥}};
  \node[n,fill=gray!8] (ani) at (0,2.1) {動物};
  \node[n,fill=gray!8] (can) at (-3.0,1.2) {カナリア};
  \node[n,fill=gray!8] (pen) at (-3.0,-1.2) {ペンギン};
  \node[n,fill=gray!8] (wing) at (3.0,1.2) {羽};
  \node[n,fill=gray!8] (fly) at (3.0,-1.2) {飛ぶ};
  \draw[arr] (bird) -- node[lab] {なかま（is a）} (ani);
  \draw[arr] (can) -- node[lab,sloped] {is a} (bird);
  \draw[arr] (pen) -- node[lab,sloped] {is a} (bird);
  \draw[arr] (bird) -- node[lab,sloped] {持つ} (wing);
  \draw[arr] (bird) -- node[lab,sloped] {できる} (fly);
\end{tikzpicture}
\end{column}
\begin{column}{0.42\textwidth}
\begin{tikzpicture}[box/.style={draw,rounded corners,text width=4.9cm,align=left,font=\small,inner sep=6pt}]
  \node[box,fill=gray!8] (a) {\textbf{点}：ことば（概念）\\\textbf{線}：関係（なかま・持つ）};
  \node[box,fill=gain!10,draw=gain,below=0.3cm of a] (b) {「カナリアに羽はある？」\\→ 線をたどって\textbf{「ある」}\\書いていないことも\textbf{推論}できる};
  \node[box,fill=boom!10,draw=boom,below=0.3cm of b] (c) {では「ペンギンは飛ぶ？」\\→ このあとの「限界」で};
\end{tikzpicture}
\end{column}
\end{columns}
\vspace{0.2em}
{\small この形は今日の後半、\textbf{ナレッジグラフ}としてもう一度出てきます}
\end{frame}

\begin{frame}{アフォーダンス：「それで何ができるか」も知識にする}
\begin{center}
\begin{tikzpicture}[o/.style={draw,rounded corners,minimum width=2.8cm,minimum height=1.1cm,align=center,font=\small},a/.style={draw,rounded corners,fill=boom!10,draw=boom,minimum width=3.4cm,minimum height=1.1cm,align=center,font=\small\bfseries},arr/.style={-{Latex[length=3mm]},very thick,gray}]
  \node[draw,rounded corners,fill=winter!12,draw=winter,minimum width=2.2cm,minimum height=3.6cm,align=center,font=\small] (r) at (0,-1.4) {\textbf{案内ロボット}};
  \foreach \y/\o/\t [count=\i] in {0/いす/座れる,-1.4/ドアの取っ手/つかんで開けられる,-2.8/エレベーターのボタン/押せる}{
    \node[a] (a\i) at (4.2,\y) {\t};
    \node[o,fill=gray!8] (o\i) at (8.3,\y) {\o};
    \draw[arr] (r.east |- a\i) -- (a\i);
    \draw[arr] (a\i) -- (o\i);
  }
\end{tikzpicture}
\end{center}
\vspace{0.4em}
\begin{itemize}\small
  \item 人はドアの取っ手を見れば、\textbf{自然に}つかんで開けられる。形が行動を「さそう」（アフォードする）
  \item 体のない AI は、それを知らない → ものに「何ができるか」を\textbf{知識として付けておく}
\end{itemize}
\end{frame}

% ============================================================
% 強みと限界
% ============================================================

\sectionframe{記号的 AI の\\[0.3em]強みと限界}

\begin{frame}{強み：判断の「理由」が読める}
\begin{center}
\begin{tikzpicture}[h/.style={draw,rounded corners,text width=3.9cm,minimum height=1.0cm,align=center,font=\bfseries},b/.style={draw,rounded corners,fill=gray!6,text width=3.9cm,minimum height=2.7cm,align=left,font=\small,inner sep=6pt,anchor=north}]
  \node[h,fill=gain!15,draw=gain] (h1) at (0,0) {透明である};
  \node[b] at ($(h1.south)+(0,-0.12)$) {・ルールが目に見える\\・まちがえたとき、\\\quad \textbf{どこがまちがいか}\\\quad すぐ分かる\\・直しやすい};
  \node[h,fill=gain!15,draw=gain] (h2) at (4.5,0) {専門知識を残せる};
  \node[b] at ($(h2.south)+(0,-0.12)$) {・ベテランの判断を\\\quad \textbf{ことばにして}残す\\・いつも同じ判断\\・知識を積み上げて\\\quad 組織の財産に};
  \node[h,fill=gain!15,draw=gain] (h3) at (9,0) {説明できる};
  \node[b] at ($(h3.south)+(0,-0.12)$) {・結論の\textbf{根拠}を示せる\\・あとから判断を\\\quad たどれる\\・\textbf{責任}の所在が\\\quad はっきりする};
\end{tikzpicture}
\end{center}
\vspace{0.4em}
\centering\large 同じ入力には、\textbf{いつも同じ答え}。動きを\textbf{保証}できる
\end{frame}

\begin{frame}{限界：世界を書ききれない}
\begin{center}
\begin{tikzpicture}[h/.style={draw,rounded corners,text width=3.9cm,minimum height=1.0cm,align=center,font=\bfseries},b/.style={draw,rounded corners,fill=gray!6,text width=3.9cm,minimum height=2.7cm,align=left,font=\small,inner sep=6pt,anchor=north}]
  \node[h,fill=boom!15,draw=boom] (h1) at (0,0) {知識を集める壁};
  \node[b] at ($(h1.south)+(0,-0.12)$) {・ルールを書くのは人\\\quad → \textbf{手間とお金}\\・ベテランも\textbf{ことばに}\\\quad \textbf{できない}「勘」\\・もれなく書けない};
  \node[h,fill=boom!15,draw=boom] (h2) at (4.5,0) {想定外に弱い};
  \node[b] at ($(h2.south)+(0,-0.12)$) {・書いていないことは\\\quad 判断できない\\・ルールどうしが\\\quad \textbf{ぶつかる}\\・例外が例外を呼ぶ};
  \node[h,fill=boom!15,draw=boom] (h3) at (9,0) {あいまいさが苦手};
  \node[b] at ($(h3.south)+(0,-0.12)$) {・「だいたい」「たぶん」\\\quad をあつかいにくい\\・写真や声のような\\\quad \textbf{ことばにならない}\\\quad 情報が苦手};
\end{tikzpicture}
\end{center}
\vspace{0.4em}
\centering 第2次 AI ブームが冬の時代に入った理由の一つが「\textbf{知識を書き尽くせない}」こと
\end{frame}

\begin{frame}{「鳥は飛ぶ」と書いたら、どうなる？}
\begin{columns}[c]
\begin{column}{0.52\textwidth}
\centering
\begin{tikzpicture}[r/.style={draw,rounded corners,text width=5.6cm,align=left,font=\small,inner sep=5pt}]
  \node[r,fill=gain!10,draw=gain] (r0) {\texttt{IF} 鳥 \texttt{THEN} 飛ぶ};
  \node[r,fill=boom!8,draw=boom,below=0.18cm of r0] (r1) {\ldots ただし\textbf{ペンギン}は飛ばない};
  \node[r,fill=boom!8,draw=boom,below=0.18cm of r1] (r2) {\ldots \textbf{ダチョウ}も飛ばない};
  \node[r,fill=boom!8,draw=boom,below=0.18cm of r2] (r3) {\ldots \textbf{ひな}はまだ飛べない};
  \node[r,fill=boom!8,draw=boom,below=0.18cm of r3] (r4) {\ldots \textbf{けが}をした鳥は飛べない};
  \node[r,fill=boom!8,draw=boom,below=0.18cm of r4] (r5) {\ldots \textbf{かごの中}の鳥は\ldots？};
  \node[font=\Large\bfseries,boom,below=0.05cm of r5] {\ldots};
\end{tikzpicture}
\end{column}
\begin{column}{0.46\textwidth}
\begin{tikzpicture}[box/.style={draw,rounded corners,text width=5.2cm,align=left,font=\small,inner sep=6pt}]
  \node[box,fill=gray!8] (a) {人なら「ふつうは飛ぶ」で\\\textbf{なんとなく}すませている};
  \node[box,fill=boom!10,draw=boom,below=0.3cm of a] (b) {ルールにすると、\textbf{例外}を\\全部書かないといけない\\→ ルールが増え、\textbf{ぶつかり}、\\\quad 直すのも大変に};
  \node[box,fill=winter!10,draw=winter,below=0.3cm of b] (c) {\textbf{1 万人に 1 人}の場合まで\\ルールを書くか？\\→ \textbf{公平性}の問題につながる};
\end{tikzpicture}
\end{column}
\end{columns}
\end{frame}

% ============================================================
% ビジネス応用と倫理
% ============================================================

\sectionframe{記号的 AI の\\[0.3em]ビジネス応用と倫理}

\begin{frame}{エキスパートシステム：専門家の判断をルールに}
\begin{center}
\begin{tikzpicture}[h/.style={draw,rounded corners,text width=3.9cm,minimum height=1.0cm,align=center,font=\bfseries},b/.style={draw,rounded corners,fill=gray!6,text width=3.9cm,minimum height=2.5cm,align=left,font=\small,inner sep=6pt,anchor=north}]
  \node[h,fill=winter!15,draw=winter] (h1) at (0,0) {金融};
  \node[b] at ($(h1.south)+(0,-0.12)$) {・融資の審査\\・不正な取引の検知\\・規制のチェック\\\quad （コンプライアンス）};
  \node[h,fill=gain!15,draw=gain] (h2) at (4.5,0) {医療};
  \node[b] at ($(h2.south)+(0,-0.12)$) {・\textbf{MYCIN}：感染症の\\\quad 診断と薬の提案\\\quad （1970年代・米国）\\・薬の飲み合わせの確認};
  \node[h,fill=boom!15,draw=boom] (h3) at (9,0) {流通};
  \node[b] at ($(h3.south)+(0,-0.12)$) {・発注のルール\\・季節の変化への対応\\・倉庫作業の自動化};
\end{tikzpicture}
\end{center}
\vspace{0.3em}
\begin{itemize}\small
  \item 1980〜90年代に広く使われた。\textbf{判断の基準がはっきりしている分野}で成果
  \item MYCIN は数百のルールで、専門医に近い水準と評価された。ただし\textbf{実際の診療には使われなかった}
\end{itemize}
{\scriptsize 出典：MYCIN はスタンフォード大学の報告（1982年）などをまとめた資料による}
\end{frame}

\begin{frame}{ルールの例と、ビジネスでの価値}
\begin{columns}[c]
\begin{column}{0.55\textwidth}
\begin{tikzpicture}[c/.style={draw,rounded corners,text width=6.4cm,align=left,font=\small,inner sep=6pt,fill=gray!6}]
  \node[c] (a) {\textbf{融資審査}\\[0.2em]\texttt{IF} 自己資本比率が業界平均より\\\quad 20 ポイント以上低い\\\texttt{THEN} リスク「高」・追加の担保を検討};
  \node[c,below=0.3cm of a] (b) {\textbf{在庫管理}\\[0.2em]\texttt{IF} 予報が「雨」\texttt{AND} 在庫がいつもの\\\quad 売上の 1.5 倍より少ない\\\texttt{THEN} いつもの 2 倍を発注・優先「高」};
\end{tikzpicture}
\end{column}
\begin{column}{0.43\textwidth}
\begin{tikzpicture}[box/.style={draw,rounded corners,text width=4.9cm,align=left,font=\small,inner sep=5pt,fill=gain!8,draw=gain}]
  \node[box] (x) {\textbf{判断の一貫性}\\担当者でぶれない};
  \node[box,below=0.2cm of x] (y) {\textbf{説明できる}\\トラブルのとき根拠を示せる};
  \node[box,below=0.2cm of y] (z) {\textbf{知識の共有}\\専門家の知識が組織の財産に};
\end{tikzpicture}
\end{column}
\end{columns}
\vspace{0.5em}
\centering\small 例は説明のために単純にしたもの（実在の基準ではない）。数字の決め方そのものが\textbf{人の判断}
\end{frame}

\thinkframe{身近なサービスやアプリで\\[0.3em]「これは\textbf{ルールで動いていそう}」\\[0.3em]と思うものを1つ\\[0.6em]{\large そう思った理由も}}

\begin{frame}{ルールは、今も身近で動いている}
\begin{center}
\begin{tikzpicture}[box/.style={draw,rounded corners,text width=3.6cm,minimum height=1.9cm,align=center,font=\small,fill=gray!6}]
  \node[box] (a) at (0,0) {\textbf{IC カードの}\\\textbf{オートチャージ}\\{\footnotesize 残高が一定より少ない\\→ チャージ}};
  \node[box] (b) at (4.2,0) {\textbf{メールの振り分け}\\{\footnotesize この差出人から\\→ このフォルダへ}};
  \node[box] (c) at (8.4,0) {\textbf{エアコン}\\{\footnotesize 室温が設定より高い\\→ 冷やす}};
  \node[box] (d) at (0,-2.3) {\textbf{ATM・ネットバンク}\\{\footnotesize 1日の上限をこえる\\→ 止める}};
  \node[box] (e) at (4.2,-2.3) {\textbf{予約システム}\\{\footnotesize 定員に達した\\→ 受け付けない}};
  \node[box,fill=winter!10,draw=winter] (f) at (8.4,-2.3) {\textbf{迷惑メールの判定}\\{\footnotesize 今は\textbf{ルール}と\\\textbf{機械学習}の組み合わせ}};
\end{tikzpicture}
\end{center}
\vspace{0.4em}
\centering「これを AI と呼ぶか」は微妙。でも\textbf{考え方は記号的 AI から}来ている
\end{frame}

\begin{frame}{倫理から見る：利点と課題}
\begin{columns}[t]
\begin{column}{0.48\textwidth}
\begin{tikzpicture}[h/.style={draw,rounded corners,text width=5.5cm,minimum height=0.9cm,align=center,font=\bfseries},b/.style={draw,rounded corners,fill=gray!6,text width=5.5cm,align=left,font=\small,inner sep=6pt}]
  \node[h,fill=gain!15,draw=gain] (h) {利点};
  \node[b,below=0.12cm of h] {・\textbf{説明できる}：判断の過程が見える\\・\textbf{たどれる}：だれが、いつ\\\quad ルールを変えたか記録が残る\\・\textbf{制御できる}：直接直せる。\\\quad 動きを保証できる};
\end{tikzpicture}
\end{column}
\begin{column}{0.48\textwidth}
\begin{tikzpicture}[h/.style={draw,rounded corners,text width=5.5cm,minimum height=0.9cm,align=center,font=\bfseries},b/.style={draw,rounded corners,fill=gray!6,text width=5.5cm,align=left,font=\small,inner sep=6pt}]
  \node[h,fill=boom!15,draw=boom] (h) {課題};
  \node[b,below=0.12cm of h] {・\textbf{書いた人の偏見}がルールに入る\\・ひとつのルールで、\textbf{だれにも}\\\quad 公平にするのはむずかしい\\・社会が変わると\textbf{古くなる}。\\\quad 直すほど手間が増える};
\end{tikzpicture}
\end{column}
\end{columns}
\vspace{0.8em}
\centering\large 「説明できる」ことと「\textbf{正しい}」ことは別
\end{frame}

\begin{frame}{融資審査の例：だれのルールか}
\begin{center}
\begin{tikzpicture}[box/.style={draw,rounded corners,text width=5.3cm,minimum height=3.0cm,align=left,font=\small,inner sep=7pt},arr/.style={-{Latex[length=3mm]},very thick,gray}]
  \node[box,fill=boom!8,draw=boom] (a) {\textbf{起こりうる問題}\\[0.3em]・特定の地域を不利にする\\・性別で判断がかたよる\\・新しい事業モデルを\\\quad 正しく評価できない};
  \node[box,fill=gain!8,draw=gain,right=1.2cm of a] (b) {\textbf{対策の例}\\[0.3em]・定期的にかたよりを確かめる\\・\textbf{いろいろな立場}の人で検証\\・ルールを見直し続ける\\\quad 仕組みを組みこむ};
  \draw[arr] (a) -- (b);
\end{tikzpicture}
\end{center}
\vspace{0.5em}
\centering 記号的 AI なら、かたよりが\textbf{ルールとして見える}から、\textbf{直せる}\\[0.2em]
\small データから学ぶ AI のかたよりは\textbf{見えにくい}（次回）
\end{frame}

% ============================================================
% 生成 AI
% ============================================================

\sectionframe{生成 AI は\\[0.3em]記号的 AI の限界を超えたのか？}

\begin{frame}{超えたように見えるところ}
\begin{center}
\begin{tikzpicture}[r/.style={draw,rounded corners,minimum height=1.05cm,align=center,font=\small},l/.style={r,text width=3.6cm,fill=boom!10,draw=boom},g/.style={r,text width=6.2cm,fill=genai!10,draw=genai},arr/.style={-{Latex[length=2.5mm]},thick,gray}]
  \node[font=\small\bfseries] at (0,0.8) {記号的 AI の限界};
  \node[font=\small\bfseries] at (6.4,0.8) {大規模言語モデル（ChatGPT・Gemini など）};
  \foreach \y/\a/\b [count=\i] in {0/{知識を集める壁}/{人が書かず、\textbf{大量の文章}から学ぶ},-1.3/{想定外に弱い}/{はじめての質問にも、\textbf{それらしく}答える},-2.6/{あいまいさが苦手}/{「だいたい」「いい感じに」も受けとめる}}{
    \node[l] (l\i) at (0,\y) {\a};
    \node[g] (g\i) at (6.4,\y) {\b};
    \draw[arr] (l\i) -- (g\i);
  }
\end{tikzpicture}
\end{center}
\vspace{0.5em}
\centering\large 3つの壁を、一度に越えたように見える。\textbf{でも、代わりに失ったものがある}
\end{frame}

\begin{frame}{ハルシネーション：「知識表現を持たない」ことの帰結}
\begin{center}
\begin{tikzpicture}[box/.style={draw,rounded corners,minimum height=1.25cm,align=center,font=\small},arr/.style={-{Latex[length=3mm]},very thick,gray}]
  \node[font=\small\bfseries,boom!80!black,anchor=west] at (-1.2,1.2) {記号的 AI};
  \node[box,fill=gray!8,text width=2.3cm] (q1) at (0,0) {質問};
  \node[box,fill=boom!12,draw=boom,text width=3.3cm] (k1) at (4.0,0) {\textbf{知識}を調べる\\{\footnotesize 事実・ルールの表}};
  \node[box,fill=gain!10,draw=gain,text width=3.6cm] (a1) at (8.6,0) {あれば答える\\\textbf{なければ「分かりません」}};
  \draw[arr] (q1) -- (k1); \draw[arr] (k1) -- (a1);
  \node[font=\small\bfseries,genai!80!black,anchor=west] at (-1.2,-1.35) {大規模言語モデル};
  \node[box,fill=gray!8,text width=2.3cm] (q2) at (0,-2.55) {質問};
  \node[box,fill=genai!12,draw=genai,text width=3.3cm] (k2) at (4.0,-2.55) {\textbf{次に来ることば}を\\もっともらしく続ける};
  \node[box,fill=boom!10,draw=boom,text width=3.6cm] (a2) at (8.6,-2.55) {知らなくても\\\textbf{文章はできてしまう}};
  \draw[arr] (q2) -- (k2); \draw[arr] (k2) -- (a2);
\end{tikzpicture}
\end{center}
\vspace{0.5em}
\centering 「事実の表」を持たず、\textbf{照らし合わせる先}がない → \textbf{自信ありげなまちがい}\\[0.2em]
\small これを\textbf{ハルシネーション}（幻覚）とよぶ。言語モデルの作られ方は次回くわしく
\end{frame}

\begin{frame}{なぜ「分かりません」と言わないのか}
\begin{columns}[c]
\begin{column}{0.50\textwidth}
\centering
{\small\textbf{選択式のテストで、答えが分からないとき}}\\[0.6em]
\begin{tikzpicture}[box/.style={draw,rounded corners,text width=2.5cm,minimum height=2.2cm,align=center,font=\small}]
  \node[box,fill=gray!10] (a) {\textbf{白紙で出す}\\[0.3em]必ず\\\textbf{0 点}};
  \node[box,fill=boom!12,draw=boom,right=0.4cm of a] (b) {\textbf{当てずっぽう}\\[0.3em]当たれば\\\textbf{点がもらえる}};
\end{tikzpicture}\\[0.6em]
{\small 正答率だけで点をつけると、\\\textbf{当てずっぽうが得}になる}
\end{column}
\begin{column}{0.48\textwidth}
\begin{tikzpicture}[box/.style={draw,rounded corners,text width=5.3cm,align=left,font=\small,inner sep=6pt}]
  \node[box,fill=genai!10,draw=genai] (a) {\textbf{OpenAI の説明（2025年9月）}\\言語モデルの学習や評価が、「分からない」より\textbf{推測して答える}ほうを高く評価してきた};
  \node[box,fill=gain!10,draw=gain,below=0.3cm of a] (b) {\textbf{提案}\\自信ありげなまちがいを\textbf{重く減点}し、「分からない」と言えたら\textbf{部分点}};
\end{tikzpicture}
\end{column}
\end{columns}
\vspace{0.4em}
\centering\small 減らす工夫は進んでいるが、\textbf{ゼロにはならない}\\[0.1em]{\scriptsize 出典：OpenAI「Why language models hallucinate」（2025年9月）}
\end{frame}

\begin{frame}{ビジネスで実際に起きたこと}
\begin{center}
\begin{tikzpicture}[h/.style={draw,rounded corners,text width=5.8cm,minimum height=0.9cm,align=center,font=\small\bfseries},b/.style={draw,rounded corners,fill=gray!6,text width=5.8cm,minimum height=2.7cm,align=left,font=\small,inner sep=6pt,anchor=north}]
  \node[h,fill=boom!15,draw=boom] (h1) at (0,0) {存在しない判例（米国・2023年）};
  \node[b] at ($(h1.south)+(0,-0.12)$) {・弁護士が ChatGPT で調べた判例を\\\quad 裁判所に出した\\・その判例は\textbf{実在しなかった}\\・弁護士らに\textbf{5,000 ドルの制裁金}};
  \node[h,fill=winter!15,draw=winter] (h2) at (6.5,0) {まちがった案内（カナダ・2024年）};
  \node[b] at ($(h2.south)+(0,-0.12)$) {・航空会社のチャットボットが、\\\quad 割引の条件を\textbf{まちがえて}案内\\・会社は「チャットボットの言った\\\quad ことは別」と主張\\・審判所は\textbf{会社の責任}と判断し、\\\quad 差額などの支払いを命じた};
\end{tikzpicture}
\end{center}
\vspace{0.1em}
\centering\small まちがえるのは AI でも、\textbf{責任を負うのは使った人・会社}\\[0.1em]
{\scriptsize 出典：Mata v. Avianca（米ニューヨーク南部地区連邦地裁、2023年6月）、Moffatt v. Air Canada（BC 州民事審判所、2024年2月）の報道。カナダの件のチャットボットが生成 AI だったかは公表されていない}
\end{frame}

\begin{frame}{RAG：調べてから答える}
\begin{center}
\begin{tikzpicture}[box/.style={draw,rounded corners,minimum height=1.6cm,align=center,font=\small},arr/.style={-{Latex[length=3mm]},very thick,gray}]
  \node[box,fill=gray!8,text width=2.3cm] (q) {\textbf{質問}\\「返品の\\期限は？」};
  \node[box,fill=boom!12,draw=boom,text width=2.9cm,right=0.55cm of q] (s) {\textbf{① 検索}\\社内の規程・\\マニュアルから\\関係する所を探す};
  \node[box,fill=genai!12,draw=genai,text width=2.9cm,right=0.55cm of s] (g) {\textbf{② 生成}\\見つけた文章を\\\textbf{根拠として}\\読んで答える};
  \node[box,fill=gain!12,draw=gain,text width=2.6cm,right=0.55cm of g] (a) {\textbf{答え}\\＋\\\textbf{根拠の場所}};
  \draw[arr] (q) -- (s); \draw[arr] (s) -- (g); \draw[arr] (g) -- (a);
\end{tikzpicture}
\end{center}
\vspace{0.5em}
\begin{itemize}\small
  \item RAG：検索拡張生成（Retrieval-Augmented Generation）。2020年に研究として発表され、\textbf{企業の生成 AI の定番}に
  \item 照らし合わせる先を\textbf{外から渡す}＝記号的 AI の「知識を調べる」を取りもどす
  \item ただし、\textbf{検索で外れた文書}を渡せばまちがう。根拠の場所を\textbf{人が確かめられる}ことが大事
\end{itemize}
{\scriptsize 出典：Lewis ほか「Retrieval-Augmented Generation for Knowledge-Intensive NLP Tasks」（NeurIPS 2020）}
\end{frame}

\begin{frame}{ナレッジグラフ：セマンティック・ネットワークの再来}
\begin{columns}[c]
\begin{column}{0.50\textwidth}
\centering
\begin{tikzpicture}[n/.style={draw,rounded corners,align=center,font=\small,minimum height=0.75cm,inner sep=4pt,fill=gray!8},arr/.style={-{Latex[length=2.5mm]},thick},lab/.style={font=\scriptsize,fill=white,inner sep=1pt}]
  \node[n,fill=boom!15,draw=boom] (a) at (0,0) {\textbf{富士山}};
  \node[n] (b) at (-2.4,1.6) {静岡県・山梨県};
  \node[n] (c) at (2.4,1.6) {山};
  \node[n] (d) at (-2.4,-1.6) {3,776 m};
  \node[n] (e) at (2.4,-1.6) {世界文化遺産};
  \draw[arr] (a) -- node[lab,sloped] {所在地} (b);
  \draw[arr] (a) -- node[lab,sloped] {なかま} (c);
  \draw[arr] (a) -- node[lab,sloped] {高さ} (d);
  \draw[arr] (a) -- node[lab,sloped] {登録} (e);
\end{tikzpicture}\\[0.3em]
{\scriptsize 形を示すための例}
\end{column}
\begin{column}{0.48\textwidth}
\begin{tikzpicture}[box/.style={draw,rounded corners,text width=5.3cm,align=left,font=\small,inner sep=6pt}]
  \node[box,fill=winter!10,draw=winter] (a) {\textbf{Google（2012年）}\\検索に「ナレッジグラフ」。\\「\textbf{文字列ではなく、もの}を」\\5 億をこえるもの、35 億をこえる事実と関係};
  \node[box,fill=genai!10,draw=genai,below=0.3cm of a] (b) {\textbf{生成 AI と組み合わせる}\\RAG の「調べる先」をグラフにする（Microsoft の GraphRAG、2024年など）\\→ 関係を\textbf{たどれる}・\textbf{確かめられる}};
\end{tikzpicture}
\end{column}
\end{columns}
\vspace{0.3em}
{\scriptsize 出典：Google 公式ブログ（2012年5月16日）、Microsoft Research（2024年）}
\end{frame}

\begin{frame}{記号的 AI の再評価：役割分担}
\begin{center}
\begin{tikzpicture}[box/.style={draw,rounded corners,text width=4.6cm,minimum height=3.3cm,align=left,font=\small,inner sep=7pt},arr/.style={{Latex[length=3mm]}-{Latex[length=3mm]},very thick,gray}]
  \node[box,fill=genai!10,draw=genai] (g) {\centering\textbf{\large 大規模言語モデル}\par\vspace{0.4em}
  ・ことばを受けとめる\\・まとめる・言いかえる\\・あいまいな質問に強い\\\textbf{苦手}：事実の保証};
  \node[box,fill=boom!10,draw=boom,right=1.6cm of g] (s) {\centering\textbf{\large 記号（知識・ルール）}\par\vspace{0.4em}
  ・事実の表・関係\\・守るべきルール\\・根拠をたどれる\\\textbf{苦手}：書ききれない};
  \draw[arr] (g) -- node[above,font=\small\bfseries] {組み合わせる} (s);
\end{tikzpicture}
\end{center}
\vspace{0.5em}
\centering{\large 古い技術が\textbf{消えた}のではなく、\textbf{弱点を補う相手}として戻ってきた}
\end{frame}

\thinkframe{生成 AI の答えを\\[0.3em]\textbf{そのまま使ってはいけない}場面は？\\[0.5em]{\large 1つあげて、\textbf{どう確かめるか}も}}

% ============================================================
% まとめ・演習
% ============================================================

\begin{frame}{まとめ}
\begin{center}
\begin{tikzpicture}[box/.style={draw,rounded corners,text width=3.9cm,minimum height=3.0cm,align=center,font=\small}]
  \node[box,fill=gray!8] (a) {\textbf{二大潮流}\\[0.5em]記号主義＝\\人が\textbf{知識を書く}\\コネクショニズム＝\\データから\textbf{学ぶ}};
  \node[box,fill=gain!8,right=0.35cm of a] (b) {\textbf{記号的 AI}\\[0.5em]ルール・分類・つながり\\\textbf{説明できる}\\でも\\\textbf{書ききれない}};
  \node[box,fill=genai!8,right=0.35cm of b] (c) {\textbf{生成 AI}\\[0.5em]知識表現を持たない\\→ ハルシネーション\\RAG・ナレッジグラフで\\\textbf{記号が戻ってきた}};
\end{tikzpicture}\\[0.8em]
{\large 「\textbf{何を根拠に}答えているのか」を問う}
\end{center}
\end{frame}

\begin{frame}{演習課題：身近な「ルールで動く仕組み」を分析する}
身近なサービス・機械・業務から、\textbf{ルールで動いていそうな仕組み}を1つ選び、次の3点について論じてください。
\begin{enumerate}
  \item その仕組みのルールを、\texttt{IF}〜\texttt{THEN}〜の形で\textbf{3つ以上}書き出す（自分の推測でよい）
  \item そのルールの\textbf{強み}（説明できる・動きを予測できる・制御できる など）と、\textbf{書ききれない例外}の例
  \item もしその仕組みを\textbf{生成 AI} に任せたら、何がよくなり、何が失われそうか。どう\textbf{確かめれば}安全に使えるか
\end{enumerate}
\vspace{0.2em}
\begin{itemize}\small
  \item 文字数の目安：\textbf{800〜2000字}。調べた内容には情報源を書く
  \item 生成 AI の利用可。ただし\textbf{自分のことばで説明できること}
\end{itemize}
\end{frame}

% ============================================================
% 付録
% ============================================================

\sectionframe{付録}

\begin{frame}{付録1：記号とルールのことば}
\footnotesize
\begin{columns}[t]
\begin{column}{0.48\textwidth}
\begin{itemize}\footnotesize
  \item \textbf{記号主義}：人の知識をことば（記号）とルールで書き表し、それを使って推論する AI の考え方。記号的 AI ともいう
  \item \textbf{コネクショニズム}：神経細胞のつながりをまねた仕組みに、データから学ばせる考え方
  \item \textbf{知識表現}：世界のことがらを、コンピュータが処理できる形（ルール・分類・関係）で書き表すこと
  \item \textbf{オントロジー}：ものごとの分類と、その上下の関係（何が何の仲間か）の体系
  \item \textbf{セマンティック・ネットワーク}：ことば（概念）を点、関係を線で結んだ網の目の表し方
\end{itemize}
\end{column}
\begin{column}{0.48\textwidth}
\begin{itemize}\footnotesize
  \item \textbf{アフォーダンス}：ものの形が、人や動物にある行動をさそうこと（取っ手は「つかむ」をさそう）。AI には「それで何ができるか」を知識として付ける
  \item \textbf{エキスパートシステム}：専門家の判断をルールとして集め、同じ判断をコンピュータにさせる仕組み。第2次 AI ブームの中心
  \item \textbf{ハルシネーション}：生成 AI が、事実でない内容をもっともらしく答えてしまうこと
  \item \textbf{RAG}（検索拡張生成）：質問に関係する文書を検索して生成 AI に渡し、それを根拠に答えさせる方法
  \item \textbf{ナレッジグラフ}：もの・人・場所などと、そのあいだの関係を大量につないだデータ
\end{itemize}
\end{column}
\end{columns}
\end{frame}

\begin{frame}{付録2：数字と出典}
\footnotesize
\begin{itemize}
  \item \textbf{MYCIN}：1970年代にスタンフォード大学で開発された、血液の感染症の診断と抗生物質の提案を行うエキスパートシステム。ルールは約500〜600（資料により異なる）。評価では専門医に近い水準とされたが、実際の診療には使われなかった
  \item \textbf{存在しない判例}：Mata v. Avianca（米ニューヨーク南部地区連邦地裁、2023年6月22日）。ChatGPT が示した実在しない判例を書面に引用した弁護士らに、5,000ドルの制裁金。米国法曹協会（ABA）などの報道
  \item \textbf{チャットボットの案内}：Moffatt v. Air Canada（ブリティッシュコロンビア州 民事審判所、2024年2月）。チャットボットが忌引き割引の条件を誤って案内し、会社の責任とされた。ABA などの報道。このチャットボットが生成 AI だったかは公表されていない
  \item \textbf{ハルシネーションの理由}：OpenAI「Why language models hallucinate」（2025年9月）。\url{https://openai.com/index/why-language-models-hallucinate/}
  \item \textbf{RAG}：Lewis ほか（NeurIPS 2020）。\url{https://arxiv.org/abs/2005.11401}
  \item \textbf{Google のナレッジグラフ}：2012年5月16日発表。5億をこえるもの、35億をこえる事実と関係（発表時）。Google 公式ブログ「Introducing the Knowledge Graph: things, not strings」
  \item \textbf{GraphRAG}：Microsoft Research（2024年）。\url{https://arxiv.org/abs/2404.16130}
\end{itemize}
\end{frame}

\begin{frame}{付録3：生成 AI と記号的 AI — できること・気をつけること}
\footnotesize
\begin{itemize}
  \item \textbf{大規模言語モデルが得意}：ことばを受けとめる、まとめる、言いかえる、あいまいな質問に答える。人がルールを書かなくても、大量の文章から学んでいる
  \item \textbf{大規模言語モデルが苦手}：答えが\textbf{事実であること}の保証。事実の表（知識表現）を持たず、照らし合わせずに「もっともらしい続き」を作るので、知らないことも文章になってしまう
  \item \textbf{記号的 AI が得意}：決まったルールを\textbf{必ず}守る、同じ入力に同じ答えを返す、判断の根拠をたどる。法令・社内規程・料金表のように「正解が決まっている」もの
  \item \textbf{組み合わせ方}
  \begin{itemize}\footnotesize
    \item RAG：社内文書などを検索して根拠として渡す。根拠の場所を示し、人が確かめられるようにする
    \item ナレッジグラフ：事実と関係を整理したデータを「調べる先」にする
    \item ルールで確かめる：生成 AI の出力を、守るべきルール（金額の上限・禁止事項など）で点検してから使う
  \end{itemize}
  \item \textbf{気をつけること}：RAG でも、検索で外れた文書や古い文書を渡せばまちがう。まちがえるのは AI でも、\textbf{責任を負うのは使った人・会社}
\end{itemize}
\end{frame}

\end{document}
